% ============================================================================
% PAQUETES - Tesis UTEM
% ============================================================================
% Este archivo contiene todos los paquetes necesarios para el documento.
% No es necesario modificarlo a menos que necesites agregar funcionalidad.
% ============================================================================

% --- Codificación y lenguaje ---
\usepackage{cmap}                    % Permite copiar-pegar texto desde el PDF sin guiones ni saltos de línea feos
\usepackage[utf8]{inputenc}          % Permite usar tildes y ñ directamente
\usepackage[T1]{fontenc}             % Codificación de fuentes para español
\usepackage[spanish,es-tabla,es-noshorthands]{babel} % Idioma español sin shorthands problemáticos
\usepackage{microtype}               % Mejora la micro-tipografía (reduce el número de palabras cortadas al final)

% --- Fuente Arial (equivalente: Helvetica) ---
\usepackage{helvet}                  % Carga la fuente Helvetica (muy similar a Arial)
\renewcommand{\familydefault}{\sfdefault} % Establece sans-serif como fuente predeterminada

% --- Geometría de página ---
\usepackage[
    letterpaper,                     % Tamaño carta
    top=4cm,                         % Margen superior: 4 cm
    left=4cm,                        % Margen izquierdo: 4 cm
    right=2.5cm,                     % Margen derecho: 2.5 cm
    bottom=2.5cm                     % Margen inferior: 2.5 cm
]{geometry}

% --- Espaciado ---
\usepackage{setspace}                % Control de interlineado
% NOTA: No se usa el paquete parskip porque interfiere con la sangría de párrafo.
% El espacio entre párrafos se configura manualmente en formato.tex.

% --- Encabezados y pies de página ---
\usepackage{fancyhdr}                % Personalización de encabezados/pies

% --- Títulos de capítulos y secciones ---
\usepackage{titlesec}                % Personalización de títulos

% --- Gráficos e imágenes ---
\usepackage{graphicx}                % Insertar imágenes
\usepackage{float}                   % Mejor control de posición de figuras/tablas
\usepackage[font=small,labelfont=bf]{caption} % Formato de pies de figura/tabla
\usepackage{chngcntr}                % Para cambiar numeración (evitar por capítulo)

% --- Tablas ---
\usepackage{array}                   % Mejoras para tablas
\usepackage{longtable}               % Tablas que ocupan más de una página
\usepackage{booktabs}                % Líneas profesionales en tablas (\toprule, \midrule, etc.)
\usepackage{multirow}                % Celdas que abarcan varias filas

% --- Comillas tipográficas (debe cargarse ANTES de biblatex) ---
\usepackage{csquotes}                % Manejo correcto de comillas en español

% --- Bibliografía APA ---
\usepackage[
    backend=biber,                   % Motor de bibliografía
    style=apa,                       % Estilo APA
    sortcites=true,                  % Ordenar citas
    sorting=nyt,                     % Ordenar por nombre, año, título
    language=spanish                 % Idioma español para la bibliografía
]{biblatex}
\addbibresource{finales/bibliografia.bib} % Archivo de bibliografía

% Permitir quiebre de línea en URLs largas de la bibliografía
\setcounter{biburllcpenalty}{7000}
\setcounter{biburlucpenalty}{8000}

% --- Hipervínculos ---
% Configuramos el paquete url nativo para que permita saltar de línea en los guiones (-)
\usepackage{url}
\expandafter\def\expandafter\UrlBreaks\expandafter{\UrlBreaks\do\-}

\usepackage[
    colorlinks=true,                 % Enlaces con color (sin recuadros)
    linkcolor=black,                 % Color de enlaces internos
    citecolor=black,                 % Color de citas
    urlcolor=blue,                   % Color de URLs
    bookmarks=true,                  % Crear marcadores en el PDF
    pdfauthor={},                    % Se llenará después
    pdftitle={}                      % Se llenará después
]{hyperref}
\usepackage{bookmark}                % Mejores marcadores en el PDF

% --- Notas al pie ---
\usepackage[bottom]{footmisc}        % Notas al pie siempre abajo

% --- Sangría del primer párrafo ---
\usepackage{indentfirst}             % Sangría también en el primer párrafo de cada sección

% --- Matemáticas (para símbolos como \square) ---
\usepackage{amssymb}                 % Símbolos matemáticos adicionales

% --- Otros ---
\usepackage{enumitem}                % Personalización de listas
\usepackage{pdfpages}                % Incluir páginas PDF externas (útil para anexos)
\usepackage{xcolor}                  % Colores (por si se necesitan)
\usepackage{emptypage}               % Páginas en blanco sin encabezado/pie

% --- TikZ (Diagramas Arquitectónicos C4, E-R y Flujos) ---
\usepackage{tikz}
\usetikzlibrary{shapes.geometric, arrows, positioning, fit, calc, babel}
